\documentclass[11pt]{article}
\usepackage[margin=25mm]{geometry}
\usepackage{amsmath,amssymb}
\usepackage[T1]{fontenc}
\usepackage[hidelinks]{hyperref}
\setlength{\parindent}{0pt}
\setlength{\parskip}{5pt}
\setlength{\emergencystretch}{3em}
\begin{document}
\begin{center}{\Large Every 0/1 polytope has Chvatal--Gomory rank zero}\\[6pt]
{\large Conjecture 00000003973}\\[5pt]
zbw123456 \quad 18 October 2026\end{center}

Use the standard meaning of a 0/1 polytope:
$P=\operatorname{conv}(V)$ for $V\subseteq\{0,1\}^n$.
Its integer hull is $P_I=\operatorname{conv}(P\cap\mathbb Z^n)$.
Since every vertex in $V$ is integral and $P$ is convex,
\[
P=\operatorname{conv}(V)\subseteq
\operatorname{conv}(P\cap\mathbb Z^n)\subseteq P.
\]
Therefore $P_I=P$.

The CG closure is the intersection of $P$ with all cuts
$a(x)\le\lfloor b\rfloor$, where $a$ is a real linear functional
integer-valued on $\mathbb Z^n$ and $a(y)\le b$ for all $y\in P$.
These are exactly integer-coefficient CG cuts.
The CG rank is the least nonnegative $k$ for which
$\operatorname{CG}^k(P)=P_I$. The zeroth iterate is $P$, so the equality
already holds at $k=0$. Hence every 0/1 polytope has rank zero.

In particular, no family of such polytopes, including partition
polytopes in this class, can have ranks at least $c n^2$ for a fixed
$c>0$ and all sufficiently large $n$. Taking a positive $n$ above the
threshold would give $0\ge c n^2>0$.
The upper-bound clause alone is trivial, but the quadratic-order
attainment and $\Theta(n^2)$ supremum clauses are false.
In every positive dimension the class is nonempty (a singleton integer
vertex suffices), and its set of ranks is $\{0\}$.

\textbf{Scope.} A polytope with fractional vertices merely contained in
$[0,1]^n$ is a different class. The proof uses the standard vertex-based
term \emph{0/1 polytope} present in both language versions and makes no
claim about arbitrary relaxations inside the cube.

\textbf{Lean correspondence.} Lean defines the integer lattice, real
convex hull, integer hull, CG closure, and rank as an \texttt{IsLeast}
statement. \texttt{integerHull\_eq} explicitly constructs integer coordinates
for 0/1 vertices and proves both inclusions.
\texttt{rank\_zero} and \texttt{no\_positive\_rank} prove the exact rank.
\texttt{no\_quadratic\_family} rules out every positive eventual quadratic
lower bound. No rank or integer-hull assertion is assumed.

\textbf{Reproduction.} The standalone project pins Lean 4.19.0 and Mathlib
commit:
\begin{center}\small\texttt{c44e0c8ee63ca166450922a373c7409c5d26b00b}\end{center}
Run \texttt{sh verify.sh}; the Lean source prints the main results'
axiom dependencies. No proof placeholders, native evaluator bypass, or
custom axioms are used.

\textbf{Source.} The official bilingual statement is included as
\texttt{conjecture.md}, checked at repository commit
\begin{center}\small\texttt{cce2d90cb298de9595798eb674d2c80712100532}\end{center}
\end{document}
